\documentclass[11pt,a4paper]{article}

% ─────────────────────────── Préambule ───────────────────────────
\usepackage[utf8]{inputenc}
\usepackage[T1]{fontenc}
\usepackage[french]{babel}
\usepackage{amsmath,amssymb,amsthm,bm}
\usepackage{booktabs}
\usepackage{graphicx}
\usepackage[margin=2.5cm]{geometry}
\usepackage{hyperref}
\usepackage{xcolor}
% pseudo-codes composés en environnements standard (pas de dépendance
% à algorithm/algpseudocode, absents de certaines installations)

\hypersetup{colorlinks=true, linkcolor=blue!60!black, citecolor=blue!60!black}

\newtheorem{proposition}{Proposition}
\newtheorem{remark}{Remarque}
\newtheorem{definition}{Définition}

% Notations
\newcommand{\R}{\mathbb{R}}
\newcommand{\T}[1]{\bm{\mathcal{#1}}}          % tenseur
\newcommand{\M}[1]{\bm{#1}}                     % matrice
\newcommand{\norf}[1]{\left\lVert #1 \right\rVert_F}
\newcommand{\noru}[1]{\left\lVert #1 \right\rVert_1}
\newcommand{\prox}{\operatorname{prox}}
\newcommand{\st}{\operatorname{S}}              % seuillage doux
\newcommand{\unfold}[2]{\M{#1}_{(#2)}}

\title{\textbf{Modèles linéaires de fusion hyperspectrale par décomposition de
Tucker couplée à c\oe ur parcimonieux :\\ architectures et algorithmes
d'optimisation}\\[1ex]
\large Étude comparative de cinq stratégies de résolution du même problème
inverse}

\author{Hamza BOUDA\\
\small Stage Assistant Ingénieur — Laboratoire LISIC, ULCO\\
\small Encadrement : projet fusion Sentinel-2 / Sentinel-3}

\date{Juillet 2026}

\begin{document}
\maketitle

\begin{abstract}
Nous décrivons en détail la famille de modèles \emph{linéaires} développée pour
la fusion d'images hyperspectrale/multispectrale (HSI--MSI), fondée sur une
décomposition de Tucker couplée dont le c\oe ur est contraint à la parcimonie.
Le modèle génératif est unique~; ce qui distingue les cinq architectures
présentées est la \emph{stratégie d'optimisation} : (i)~moindres carrés
alternés avec étape proximale accélérée (ALS+FISTA), (ii)~gradient adaptatif
avec opérateur proximal découplé (Adam+prox), (iii)~gradient adaptatif avec
pénalité $\ell_1$ sous-différentiée (Adam+$\ell_1$), (iv)~gradient stochastique
à inertie avec étape proximale (SGD+prox), et (v)~quasi-Newton à mémoire
limitée avec étape proximale (LBFGS+prox). Pour chaque variante nous donnons
les équations de mise à jour complètes, les garanties théoriques, le coût par
itération et le traitement exact ou approché de la contrainte de parcimonie.
Nous montrons que ces choix, à modèle strictement identique, induisent des
comportements très différents~: l'écart entre solveurs peut dépasser l'écart
entre classes de modèles, ce qui impose de séparer soigneusement
\og{}contribution du modèle\fg{} et \og{}contribution de l'optimiseur\fg{} dans
toute étude comparative.
\end{abstract}

% ══════════════════════════════════════════════════════════════════
\section{Introduction}
% ══════════════════════════════════════════════════════════════════

La fusion HSI--MSI vise à reconstruire une image de haute résolution à la fois
spatiale et spectrale, la \emph{super-resolution image} (SRI)
$\T{S}\in\R^{H\times W\times C}$, à partir de deux observations dégradées et
complémentaires~:
\begin{itemize}
  \item une image \textbf{hyperspectrale basse résolution spatiale}
  $\T{H}\in\R^{h\times w\times C}$ (riche spectralement, $C$ bandes, grille
  grossière $h\times w$ avec $H/h = W/w = r$)~;
  \item une image \textbf{multispectrale haute résolution spatiale}
  $\T{M}\in\R^{H\times W\times c}$ (fine spatialement, pauvre spectralement,
  $c\ll C$ bandes).
\end{itemize}
Ce problème est le c\oe ur méthodologique de la fusion Sentinel-2/Sentinel-3~:
Sentinel-2 joue le rôle du flux multispectral fin (jusqu'à 10\,m, 13 bandes) et
Sentinel-3/OLCI celui du flux hyperspectral grossier (300\,m, 21 bandes).

Le problème est un \emph{problème inverse mal posé}~: le nombre d'inconnues
($H\!\cdot\!W\!\cdot\!C$) excède largement le nombre d'observations. Toute
méthode doit donc injecter un \emph{a priori} de structure. La famille de
modèles étudiée ici adopte l'a priori multilinéaire de \textsc{Tucker}~: la
scène vit dans le produit de trois sous-espaces de faible dimension (spatial
ligne, spatial colonne, spectral), couplés par un c\oe ur creux qui encode les
interactions.

\paragraph{Positionnement.} Ce cadre est celui des méthodes tensorielles de la
littérature — STEREO (CPD couplée), SCOTT (Tucker couplée par HOSVD), et leurs
variantes parcimonieuses G-STEREO/G-SCOTT. Notre contribution dans ce rapport
n'est pas le modèle lui-même mais son étude \emph{à solveur variable}~: nous
résolvons \textbf{le même} programme d'optimisation par cinq algorithmes de
natures différentes et nous explicitons ce que chaque choix change,
mathématiquement et pratiquement.

% ══════════════════════════════════════════════════════════════════
\section{Notations et rappels d'algèbre multilinéaire}
% ══════════════════════════════════════════════════════════════════

Les tenseurs sont notés en calligraphie grasse ($\T{S}$), les matrices en gras
($\M{D}$), les scalaires en italique. Pour $\T{X}\in\R^{I_1\times I_2\times
I_3}$~:

\begin{definition}[Dépliement]
Le dépliement selon le mode $n$, noté $\unfold{X}{n}$, réarrange les fibres de
mode $n$ en colonnes~: $\unfold{X}{1}\in\R^{I_1\times I_2 I_3}$,
$\unfold{X}{2}\in\R^{I_2\times I_1 I_3}$, $\unfold{X}{3}\in\R^{I_3\times I_1
I_2}$.
\end{definition}

\begin{definition}[Produit de mode]
Le produit de mode $n$ entre $\T{X}$ et $\M{U}\in\R^{J\times I_n}$, noté
$\T{X}\times_n \M{U}$, contracte la dimension $I_n$~:
$(\T{X}\times_1\M{U})_{j i_2 i_3} = \sum_{i_1} U_{j i_1} X_{i_1 i_2 i_3}$,
soit en dépliement $\unfold{(\T{X}\times_n\M{U})}{n} = \M{U}\,\unfold{X}{n}$.
\end{definition}

\begin{definition}[Décomposition de Tucker]
Une décomposition de Tucker de rangs multilinéaires $(R_1,R_2,R_3)$ s'écrit
\begin{equation}
  \T{S} \;=\; \T{G}\times_1 \M{D}_1 \times_2 \M{D}_2 \times_3 \M{D}_3,
  \qquad
  \T{G}\in\R^{R_1\times R_2\times R_3},\;
  \M{D}_n\in\R^{I_n\times R_n},
  \label{eq:tucker}
\end{equation}
où $\T{G}$ est le \emph{c\oe ur} et les $\M{D}_n$ les \emph{dictionnaires} (ou
facteurs) de mode. En coordonnées~:
$S_{ijk}=\sum_{a=1}^{R_1}\sum_{b=1}^{R_2}\sum_{c=1}^{R_3}
G_{abc}\,(D_1)_{ia}(D_2)_{jb}(D_3)_{kc}$.
\end{definition}

Le nombre de degrés de liberté passe de $HWC$ à
$R_1R_2R_3 + HR_1 + WR_2 + CR_3$ — pour Pavia ($608\times336\times103$) et des
rangs $(30,30,10)$, c'est une compression d'un facteur $\sim 700$. C'est cet
écrasement dimensionnel qui rend le problème inverse identifiable.

% ══════════════════════════════════════════════════════════════════
\section{Modèle d'observation}
\label{sec:observation}
% ══════════════════════════════════════════════════════════════════

\subsection{Dégradation spatiale (voie hyperspectrale)}

L'acquisition basse résolution est modélisée par un flou de réponse
impulsionnelle (PSF) suivi d'une décimation de facteur $r$, \emph{séparable}
par lignes et colonnes~:
\begin{equation}
  \T{H} \;=\; \T{S}\times_1 \M{B}_h \times_2 \M{B}_w \;+\; \T{E}_H,
  \qquad \M{B}_h\in\R^{h\times H},\; \M{B}_w\in\R^{w\times W},
  \label{eq:degH}
\end{equation}
où $\T{E}_H$ est un bruit blanc gaussien. La matrice $\M{B}_h$ est construite
depuis la coupe centrale 1-D $p\in\R^{2m+1}$ du noyau PSF 2-D fourni
(normalisée à somme unité), par
\begin{equation}
  (\M{B}_h)_{i,\,j} \;=\; \frac{p_{\,j-ir+m}}
  {\sum_{l} p_{\,l-ir+m}}
  \quad\text{si } |j-ir|\le m,\qquad 0\text{ sinon},
  \label{eq:psfmat}
\end{equation}
c'est-à-dire une convolution-décimation à bords renormalisés (chaque ligne de
$\M{B}_h$ somme à 1, ce qui préserve la radiométrie moyenne).

\subsection{Dégradation spectrale (voie multispectrale)}

L'intégration des bandes fines dans les bandes larges du capteur MSI est
modélisée par la matrice de réponses spectrales (SRF) $\M{R}\in\R^{c\times
C}$~:
\begin{equation}
  \T{M} \;=\; \T{S}\times_3 \M{R} \;+\; \T{E}_M .
  \label{eq:degM}
\end{equation}

\subsection{Normalisation et cohérence d'échelle}

Chaque observation est normalisée par son percentile 99~:
$\widetilde{\T{H}}=\T{H}/s_H$, $\widetilde{\T{M}}=\T{M}/s_M$ avec
$s_X=\mathrm{perc}_{99}(\T{X})$. La SRI est estimée dans l'échelle de
$\T{H}$~; la prédiction MSI doit donc être recalée par le facteur
$\rho = s_H/s_M$~:
\begin{equation}
  \widehat{\T{M}} \;=\; \rho\,\bigl(\T{S}\times_3\M{R}\bigr).
  \label{eq:scale}
\end{equation}
Cette correction, absente des implémentations historiques, supprime un biais
systématique entre les deux termes d'attache aux données.

% ══════════════════════════════════════════════════════════════════
\section{Le modèle commun : Tucker couplé à c\oe ur parcimonieux}
\label{sec:model}
% ══════════════════════════════════════════════════════════════════

Toutes les architectures de ce rapport estiment le \emph{même} objet~: le
quadruplet $(\T{G},\M{D}_1,\M{D}_2,\M{D}_3)$ minimisant la fonction de coût
couplée
\begin{equation}
\boxed{\;
  \mathcal{L}(\T{G},\M{D}_1,\M{D}_2,\M{D}_3)
  \;=\;
  \underbrace{\lambda_H\,
  \norf{\widetilde{\T{H}} - \T{G}\times_1\M{B}_h\M{D}_1
  \times_2\M{B}_w\M{D}_2\times_3\M{D}_3}^2}_{\text{attache HSI}}
  \;+\;
  \underbrace{\lambda_M\,
  \norf{\widetilde{\T{M}} - \rho\,\T{G}\times_1\M{D}_1
  \times_2\M{D}_2\times_3\M{R}\M{D}_3}^2}_{\text{attache MSI}}
  \;+\;
  \underbrace{\beta\,\noru{\T{G}}}_{\text{parcimonie}}
\;}
\label{eq:loss}
\end{equation}
avec $\noru{\T{G}}=\sum_{abc}|G_{abc}|$ et, dans toutes nos expériences,
$\lambda_H=\lambda_M=1$.

\subsection{Rôle de chaque terme}

\textbf{Attache HSI.} La reconstruction $\T{S}$ dégradée spatialement
(éq.~\ref{eq:degH}) doit expliquer l'observation hyperspectrale~: ce terme
ancre les \emph{spectres}. Comme $\M{B}_h\M{D}_1$ compose le dictionnaire
spatial avec le flou, l'information spatiale fine de $\M{D}_1$ n'y est pas
contrainte — elle l'est par le second terme.

\textbf{Attache MSI.} La reconstruction projetée spectralement
(éq.~\ref{eq:degM}) doit expliquer l'observation multispectrale~: ce terme
ancre la \emph{géométrie fine}. La complémentarité des deux projections
($\M{B}$ agit sur les modes 1--2, $\M{R}$ sur le mode 3) est ce qui rend le
couplage identifiant~: aucune observation seule ne contraint tous les modes.

\textbf{Parcimonie du c\oe ur.} La pénalité $\ell_1$ sur $\T{G}$ est la
relaxation convexe du comptage $\ell_0$. Elle repose sur l'hypothèse physique
que \emph{peu d'interactions triples} (motif spatial ligne $\times$ motif
spatial colonne $\times$ signature spectrale) suffisent à décrire la scène~:
le c\oe ur agit comme une table de contingence creuse entre les atomes des
trois dictionnaires. Outre la régularisation statistique, la parcimonie
protège contre le sur-ajustement du bruit dans le régime très sous-déterminé.

\subsection{Ambiguïté d'échelle et normalisation des dictionnaires}
\label{sec:ambiguity}

\begin{proposition}[Invariance d'échelle]
Pour toutes matrices diagonales inversibles $\M{\Lambda}_n$, le remplacement
$\M{D}_n \leftarrow \M{D}_n\M{\Lambda}_n^{-1}$,
$\T{G} \leftarrow \T{G}\times_1\M{\Lambda}_1\times_2\M{\Lambda}_2
\times_3\M{\Lambda}_3$ laisse $\T{S}$ inchangée mais modifie $\noru{\T{G}}$.
\end{proposition}

Sans contrainte supplémentaire, l'optimiseur peut donc \emph{contourner} la
pénalité $\ell_1$ en gonflant les dictionnaires et en écrasant le c\oe ur
(ou l'inverse), rendant le seuillage incohérent d'une itération à l'autre.
Toutes nos architectures lèvent cette ambiguïté par la \textbf{normalisation
des colonnes}~:
\begin{equation}
  \M{D}_n \;\longleftarrow\; \M{D}_n\,
  \mathrm{diag}\bigl(\lVert \M{d}_{n,1}\rVert_2,\dots,
  \lVert \M{d}_{n,R_n}\rVert_2\bigr)^{-1},
  \label{eq:normcols}
\end{equation}
appliquée à chaque itération (ALS) ou à chaque passe avant (variantes
gradient, où la normalisation est \emph{incluse dans le graphe de calcul} et
donc différentiée). Toute l'énergie vit alors dans $\T{G}$, et $\beta$ a une
signification stable.

\subsection{Initialisation par SVD des observations}
\label{sec:init}

Le problème (\ref{eq:loss}) est non convexe~; le point de départ est décisif.
Toutes les architectures partagent la même initialisation \emph{spectrale
par les données}, qui n'utilise que les observations (jamais la référence)~:
\begin{align}
  \M{D}_1^{(0)} &= U_{R_1}\!\bigl(\unfold{\widetilde{M}}{1}\bigr), &
  \M{D}_2^{(0)} &= U_{R_2}\!\bigl(\unfold{\widetilde{M}}{2}\bigr), &
  \M{D}_3^{(0)} &= U_{R_3}\!\bigl(\unfold{\widetilde{H}}{3}\bigr),
  \label{eq:init-svd}
\end{align}
où $U_k(\cdot)$ désigne les $k$ premiers vecteurs singuliers gauches~: les
sous-espaces spatiaux viennent du flux le plus fin spatialement, le
sous-espace spectral du flux le plus riche spectralement. Le c\oe ur est
initialisé par projection de l'observation fine~:
\begin{equation}
  \T{G}^{(0)} \;=\;
  \widetilde{\T{M}}\times_1{\M{D}_1^{(0)}}^{\!\top}
  \times_2{\M{D}_2^{(0)}}^{\!\top}
  \times_3\bigl(\M{R}\M{D}_3^{(0)}\bigr)^{\!\top},
  \qquad
  \T{G}^{(0)} \leftarrow \T{G}^{(0)}/\norf{\T{G}^{(0)}}.
  \label{eq:init-core}
\end{equation}

% ══════════════════════════════════════════════════════════════════
\section{Architecture I --- ALS + FISTA (solveur de référence)}
\label{sec:als}
% ══════════════════════════════════════════════════════════════════

La première architecture suit la tradition des méthodes tensorielle (STEREO,
SCOTT, CTF)~: une \textbf{minimisation alternée par blocs}. On observe que
(\ref{eq:loss}), non convexe globalement, est \emph{convexe par bloc}~: à
trois blocs fixés, le problème dans le quatrième est un moindre carré
(éventuellement pénalisé $\ell_1$). L'algorithme boucle donc sur quatre
sous-problèmes résolus exactement ou quasi exactement.

\subsection{Mise à jour des dictionnaires (moindres carrés)}

Écrivons le résidu de la source $k$ en dépliement de mode $n$~: en notant
$\M{B}_k^{(n)}$ l'opérateur agissant sur le mode $n$ pour la source $k$
(p.ex. $\M{B}_h$ pour la source HSI en mode 1, $\M{I}$ pour la source MSI) et
$\M{W}_k^{(n)}$ le produit de Khatri--Rao--Kronecker des autres facteurs
dégradés,
\begin{equation}
  \unfold{Y_k}{n} \;\approx\;
  \M{B}_k^{(n)}\,\M{D}_n\;\M{A}_k^{(n)\top},
  \qquad
  \M{A}_k^{(n)} := \text{(facteurs et c\oe ur des autres modes)} .
\end{equation}
La mise à jour pratique de $\M{D}_n$ agrège toutes les sources dans un
système linéaire de type normal~:
\begin{equation}
  \Bigl[\;\sum_k \lambda_k\,
  \M{B}_k^{(n)\top}\M{B}_k^{(n)} + \varepsilon\M{I}\;\Bigr]\,
  \M{D}_n
  \;=\;
  \sum_k \lambda_k\,
  \M{B}_k^{(n)\top}\,\unfold{Y_k}{n}\,\M{A}_k^{(n)},
  \label{eq:als-D}
\end{equation}
résolu par moindres carrés régularisés ($\varepsilon=10^{-5}$), puis
normalisation des colonnes (éq.~\ref{eq:normcols}). Chaque source contribue au
système~: le dictionnaire spatial voit à la fois la version floutée (via
$\M{B}^\top\M{B}$) et la version fine (via $\M{I}$).

\subsection{Mise à jour du c\oe ur (FISTA)}

À dictionnaires fixés, le sous-problème en $\T{G}$ est un LASSO
tensoriel~: $\min_{\T{G}} f(\T{G}) + \beta\noru{\T{G}}$ avec $f$ quadratique.
Son gradient s'obtient par rétro-projection des résidus~:
\begin{equation}
  \nabla f(\T{G})
  \;=\;
  2\sum_k \lambda_k\;
  \T{R}_k \times_1 (\M{B}_k^{(1)}\M{D}_1)^{\!\top}
          \times_2 (\M{B}_k^{(2)}\M{D}_2)^{\!\top}
          \times_3 (\M{B}_k^{(3)}\M{D}_3)^{\!\top},
  \label{eq:gradG}
\end{equation}
où $\T{R}_k$ est le résidu de la source $k$. La constante de Lipschitz de
$\nabla f$ est majorée par le produit des normes spectrales~:
\begin{equation}
  L \;=\; 2\sum_k \lambda_k\,
  \prod_{n=1}^{3}\bigl\lVert \M{B}_k^{(n)}\M{D}_n\bigr\rVert_2^2 ,
  \label{eq:lipschitz}
\end{equation}
ce qui donne un pas \emph{calculé, sans réglage}. L'itération FISTA combine
étape de gradient, opérateur proximal et extrapolation de Nesterov~:
\begin{align}
  \T{G}^{(j+1)} &= \st_{\beta/L}\!\Bigl(\T{Z}^{(j)}
     - \tfrac{1}{L}\nabla f(\T{Z}^{(j)})\Bigr),
  &t_{j+1} &= \tfrac{1+\sqrt{1+4t_j^2}}{2},
  &\T{Z}^{(j+1)} &= \T{G}^{(j+1)}
     + \tfrac{t_j-1}{t_{j+1}}\bigl(\T{G}^{(j+1)}-\T{G}^{(j)}\bigr),
  \label{eq:fista}
\end{align}
où $\st_\tau(x)=\mathrm{sign}(x)\max(|x|-\tau,0)$ est le \textbf{seuillage
doux}, opérateur proximal exact de $\tau\noru{\cdot}$. FISTA garantit
$f(\T{G}^{(j)})-f^\star = O(1/j^2)$ sur ce sous-problème.

\subsection{Propriétés}

\begin{itemize}
  \item \textbf{Garanties}~: descente monotone de $\mathcal{L}$ à chaque bloc~;
  convergence vers un point critique du problème joint (au sens des méthodes
  BCD sur fonctions multi-convexes).
  \item \textbf{Parcimonie exacte}~: le prox produit de vrais zéros~;
  la sparsité de $\T{G}$ est mesurable et interprétable.
  \item \textbf{Hyperparamètres}~: quasi aucun ($\beta$, rangs, nombre
  d'itérations)~; le pas est déduit de (\ref{eq:lipschitz}).
  \item \textbf{Coût}~: chaque itération externe requiert des dépliements, des
  systèmes linéaires $R_n\times R_n$ par bloc et $O(n_{\text{FISTA}})$
  produits de mode~; c'est l'architecture la plus chère par itération
  ($\sim$265\,s par scène Pavia en pratique, CPU).
  \item \textbf{Limite}~: la résolution \emph{exacte} de chaque bloc peut
  sur-adapter le bloc courant aux erreurs des blocs figés, et la trajectoire
  alternée peut converger vers des points critiques de mauvaise qualité —
  phénomène observé expérimentalement (voir~\S\ref{sec:discussion}).
\end{itemize}

% ══════════════════════════════════════════════════════════════════
\section{Le cadre gradient unifié (architectures II--V)}
\label{sec:gradient}
% ══════════════════════════════════════════════════════════════════

Les quatre architectures suivantes abandonnent l'alternance~: tous les
paramètres $\theta=(\T{G},\M{D}_1,\M{D}_2,\M{D}_3)$ sont mis à jour
\emph{simultanément} par différentiation automatique de la partie lisse de
(\ref{eq:loss}),
\begin{equation}
  f(\theta) \;=\; \lambda_H\norf{\widetilde{\T{H}}-\widehat{\T{H}}(\theta)}^2
  + \lambda_M\norf{\widetilde{\T{M}}-\widehat{\T{M}}(\theta)}^2 ,
\end{equation}
la normalisation des colonnes (\ref{eq:normcols}) étant \emph{à l'intérieur}
du graphe de calcul~: on optimise des paramètres libres
$\M{D}_n^{\text{raw}}$ et le modèle utilise
$\M{D}_n=\M{D}_n^{\text{raw}}/\lVert\cdot\rVert_{2,\text{col}}$, si bien que
la contrainte est exactement satisfaite à chaque évaluation, sans projection
externe. Les gradients sont bornés par écrêtage de norme globale
($\lVert\nabla\rVert_2\le 10$).

La pénalité $\ell_1$, non lisse, est traitée de deux manières concurrentes qui
définissent les sous-familles \emph{prox} et \emph{$\ell_1$-pénalisée}~:

\paragraph{(a) Schéma proximal (splitting avant--arrière).}
Après chaque pas de l'optimiseur sur $f$ seul, on applique le prox au c\oe ur~:
\begin{equation}
  \T{G} \;\longleftarrow\; \st_{\tau}(\T{G}),
  \qquad \tau = \beta\,\eta_0 ,
  \label{eq:prox-step}
\end{equation}
où $\eta_0$ est le \textbf{pas nominal} (et non le pas courant du
scheduler)~: ce \emph{découplage} évite que la décroissance du taux
d'apprentissage n'affaiblisse silencieusement le seuil effectif en fin
d'entraînement. Comme en FISTA, le prox produit des \textbf{zéros exacts}.

\paragraph{(b) Pénalité sous-différentiée.}
On ajoute $\beta\noru{\T{G}}$ directement à la perte~; l'autodiff propage
$\beta\,\mathrm{sign}(\T{G})$. Le gradient pousse les coefficients vers zéro
mais, un pas de gradient étant une opération continue, \textbf{aucun
coefficient n'atteint exactement zéro}~: la parcimonie est approchée
(coefficients $O(\beta\eta)$), non structurelle.

\begin{remark}
La distinction (a)/(b) est invisible sur la fidélité lorsque $\beta$ est
petit (les deux variantes coïncident à $10^{-4}$ près dans nos mesures pour
$\beta=10^{-6}$) mais devient essentielle dès que la parcimonie est un
objectif en soi~: seul le schéma proximal fournit un support creux
exploitable et un diagnostic fiable de la santé du c\oe ur.
\end{remark}

% ══════════════════════════════════════════════════════════════════
\section{Architecture II --- Adam + prox}
\label{sec:adamprox}
% ══════════════════════════════════════════════════════════════════

L'optimiseur Adam maintient des moments d'ordre 1 et 2 par coordonnée~:
\begin{align}
  \M{m}_t &= \gamma_1\M{m}_{t-1} + (1-\gamma_1)\,\M{g}_t, &
  \M{v}_t &= \gamma_2\M{v}_{t-1} + (1-\gamma_2)\,\M{g}_t^{\odot 2}, \\
  \hat{\M{m}}_t &= \M{m}_t/(1-\gamma_1^t), &
  \hat{\M{v}}_t &= \M{v}_t/(1-\gamma_2^t), \\
  \theta_{t+1} &= \theta_t - \eta_t\,
  \frac{\hat{\M{m}}_t}{\sqrt{\hat{\M{v}}_t}+\epsilon},
  &&(\gamma_1,\gamma_2,\epsilon)=(0.9,\,0.999,\,10^{-8}),
  \label{eq:adam}
\end{align}
suivi de l'étape proximale (\ref{eq:prox-step}) sur $\T{G}$ uniquement.

\paragraph{Pourquoi Adam domine-t-il sur ce problème ?}
Le paysage de (\ref{eq:loss}) est extrêmement \textbf{mal conditionné}~: les
sensibilités de $\mathcal{L}$ aux coordonnées de $\T{G}$ (multipliées par
trois dictionnaires) et aux coordonnées des $\M{D}_n$ (multipliées par le
c\oe ur et deux dictionnaires) diffèrent de plusieurs ordres de grandeur, et
varient au cours de l'optimisation. La renormalisation par
$\sqrt{\hat{\M{v}}_t}$ agit comme un préconditionneur diagonal adaptatif qui
égalise ces échelles — c'est l'analogue stochastique du pas de Lipschitz par
bloc de l'ALS, mais \emph{par coordonnée} et sans calcul de normes
spectrales.

\paragraph{Réglage retenu (après balayage).} $\eta_0=10^{-2}$, scheduler
cosinus $\eta_t = \eta_{\min} + \tfrac12(\eta_0-\eta_{\min})
(1+\cos(\pi t/T))$, $T=10^4$ itérations, arrêt anticipé sur la perte
(patience réglable), seuil prox découplé $\tau=\beta\eta_0$.

\paragraph{Complexité.} Une itération $=$ une passe avant (deux produits de
mode complets) $+$ une rétropropagation ($\approx 2\times$ la passe avant)
$+$ $O(\#\theta)$ pour Adam~; entièrement vectorisée sur GPU
($\sim$40--90\,s pour $10^4$ itérations sur A100, scène Pavia).

% ══════════════════════════════════════════════════════════════════
\section{Architecture III --- Adam + pénalité $\ell_1$}
\label{sec:adaml1}
% ══════════════════════════════════════════════════════════════════

Identique à l'architecture II \emph{sauf} le traitement de la parcimonie~:
la perte optimisée devient $f(\theta)+\beta\noru{\T{G}}$ et aucune étape
proximale n'est appliquée. Le sous-gradient utilisé par l'autodiff est
\begin{equation}
  \partial_{\T{G}}\bigl[\beta\noru{\T{G}}\bigr]
  = \beta\,\mathrm{sign}(\T{G})
  \quad(\text{avec } \mathrm{sign}(0)=0).
\end{equation}

\paragraph{Ce qui change.}
\begin{itemize}
  \item \textbf{Pas de zéros exacts}~: au voisinage de $0$, le terme de
  fidélité et la pénalité s'équilibrent en un point $\neq 0$~; le support de
  $\T{G}$ reste plein. La \og{}parcimonie\fg{} est un rétrécissement
  ($shrinkage$), pas une sélection.
  \item \textbf{Interaction avec Adam}~: la pénalité passe par le
  préconditionneur $1/\sqrt{\hat{v}}$~; le rétrécissement effectif dépend
  donc de la variance du gradient de chaque coordonnée — un effet de bord
  connu (analogue au problème réglé par AdamW pour le weight decay)~: le
  \emph{même} $\beta$ ne produit pas la \emph{même} contraction partout.
  Le schéma proximal (架. II), appliqué hors du préconditionneur, n'a pas ce
  défaut.
  \item \textbf{Intérêt}~: simplicité maximale (une ligne), et sert de
  contrôle expérimental — l'écart II$\,$vs$\,$III isole \emph{l'effet du
  traitement de la non-lissité} à optimiseur identique.
\end{itemize}

% ══════════════════════════════════════════════════════════════════
\section{Architecture IV --- SGD à inertie + prox}
\label{sec:sgd}
% ══════════════════════════════════════════════════════════════════

Descente de gradient classique avec moment de Polyak~:
\begin{equation}
  \M{u}_t = \mu\,\M{u}_{t-1} + \M{g}_t, \qquad
  \theta_{t+1} = \theta_t - \eta\,\M{u}_t,
  \qquad \mu = 0.9,\;\; \eta = \eta_0/2,
  \label{eq:sgd}
\end{equation}
suivie de l'étape proximale (\ref{eq:prox-step}). Le pas est réduit de moitié
par rapport à Adam car, sans renormalisation par coordonnée, le pas stable est
dicté par la \emph{pire} courbure du paysage.

\paragraph{Ce qui change et ce que cela révèle.}
SGD applique le \emph{même pas à toutes les coordonnées}. Sur un problème
multilinéaire mal conditionné, cela force un compromis~: pas assez grand pour
faire progresser les directions plates (dictionnaires), pas assez petit pour
ne pas osciller dans les directions raides (c\oe ur). Expérimentalement, SGD
suit Adam pendant la phase initiale (chute rapide portée par l'init SVD) puis
\emph{décroche} dans la phase de raffinement — l'écart final ($\approx$
$-0.3$ à $-0.6$\,dB selon la scène) mesure précisément la valeur du
préconditionnement adaptatif. L'inertie ($\mu=0.9$) atténue mais ne supprime
pas ce déficit~: elle accélère dans les vallées cohérentes mais n'égalise pas
les échelles.

% ══════════════════════════════════════════════════════════════════
\section{Architecture V --- LBFGS + prox}
\label{sec:lbfgs}
% ══════════════════════════════════════════════════════════════════

L-BFGS approxime l'inverse du hessien par une histoire de $m$ paires
$(\M{s}_i,\M{y}_i)=(\theta_{i+1}-\theta_i,\;\M{g}_{i+1}-\M{g}_i)$~:
\begin{equation}
  \theta_{t+1} = \theta_t - \eta\,\M{H}_t\,\M{g}_t,
  \qquad
  \M{H}_t \approx (\nabla^2 f)^{-1}
  \text{ (récurrence à deux boucles, mémoire } m=10\text{)},
  \label{eq:lbfgs}
\end{equation}
avec recherche linéaire interne (jusqu'à 20 évaluations de $f$ par pas
externe), puis étape proximale (\ref{eq:prox-step}).

\paragraph{Ce qui change.}
\begin{itemize}
  \item \textbf{Information de courbure}~: contrairement à Adam (diagonal),
  LBFGS capte des \emph{corrélations entre coordonnées} — utile quand les
  blocs $\T{G}$ et $\M{D}_n$ interagissent fortement.
  \item \textbf{Coût par pas}~: chaque pas externe vaut jusqu'à 20 passes
  avant/arrière~; à budget temps égal, LBFGS fait $\sim$20$\times$ moins de
  pas. On le lance donc sur $\sim$300 pas externes avec $\eta=0.5$.
  \item \textbf{Friction théorique}~: la théorie BFGS suppose $f$ lisse~;
  l'étape proximale intercalée viole cette hypothèse (le prox déplace
  $\theta$ hors de la trajectoire prédite par $\M{H}_t$). Avec $\beta$ petit,
  la perturbation est négligeable~; pour $\beta$ grand, un schéma
  \emph{prox-quasi-Newton} dédié serait nécessaire.
  \item \textbf{Note d'implémentation}~: l'implémentation PyTorch de LBFGS
  aplati les gradients par \texttt{.view(-1)}, ce qui exige des gradients
  \emph{contigus} en mémoire — les gradients issus d'\texttt{einsum} ne le
  sont pas toujours~; ils sont explicitement rendus contigus dans la
  fermeture d'évaluation.
\end{itemize}

% ══════════════════════════════════════════════════════════════════
\section{Synthèse comparative}
\label{sec:discussion}
% ══════════════════════════════════════════════════════════════════

\subsection{Tableau récapitulatif}

\begin{table}[h]
\centering
\small
\begin{tabular}{@{}lccccc@{}}
\toprule
 & \textbf{I. ALS+FISTA} & \textbf{II. Adam+prox} & \textbf{III. Adam+$\ell_1$}
 & \textbf{IV. SGD+prox} & \textbf{V. LBFGS+prox} \\
\midrule
Mise à jour            & alternée, exacte & simultanée & simultanée & simultanée & simultanée \\
Pas                    & calculé (Lipschitz) & adaptatif/coord. & adaptatif/coord. & global fixe & courbure $m{=}10$ \\
Zéros exacts dans $\T{G}$ & \textbf{oui} (prox) & \textbf{oui} (prox) & non & \textbf{oui} (prox) & \textbf{oui} (prox) \\
Descente monotone      & \textbf{oui} & non & non & non & non (recherche lin.) \\
Hyperparamètres        & $\beta$, rangs & $+\,\eta_0$, scheduler & $+\,\eta_0$, scheduler & $+\,\eta$, $\mu$ & $+\,\eta$, $m$ \\
Coût/itération         & élevé & faible & faible & faible & $\sim$20$\times$ faible \\
Support GPU            & partiel & natif & natif & natif & natif \\
Extension non-linéaire & re-dérivation totale & \emph{immédiate} & \emph{immédiate} & \emph{immédiate} & délicate \\
\bottomrule
\end{tabular}
\caption{Les cinq architectures~: même modèle (éq.~\ref{eq:loss}), cinq
stratégies de résolution.}
\label{tab:comparison}
\end{table}

\subsection{Ce qui change d'une architecture à l'autre — lecture progressive}

\textbf{I $\to$ II (alternance $\to$ gradient simultané).} Le changement le
plus lourd de conséquences. L'ALS résout chaque bloc \emph{exactement} dans
le contexte figé des autres — optimal localement, mais la trajectoire
alternée peut s'enfermer dans des points critiques médiocres où chaque bloc
est optimal contre les erreurs des autres. Le gradient simultané fait des pas
\emph{imparfaits mais cohérents}~: tous les paramètres co-évoluent, et le
bruit de discrétisation aide à franchir les selles. Expérimentalement, cet
unique changement vaut plusieurs dB à modèle identique — de loin l'écart
dominant de l'étude, supérieur à ce que l'on gagne ensuite en changeant
d'optimiseur au sein de la famille gradient.

\textbf{II $\to$ III (prox $\to$ pénalité).} Ne change ni la fidélité (à
$\beta$ faible) ni la vitesse~; change la \emph{nature} de la parcimonie
(exacte vs approchée) et sa \emph{robustesse} (seuil hors préconditionneur vs
rétrécissement modulé par $\hat{v}$). C'est un contrôle expérimental plus
qu'une alternative.

\textbf{II $\to$ IV (adaptatif $\to$ pas global).} Supprime le
préconditionnement diagonal~; mesure directement le coût du mauvais
conditionnement multilinéaire. Le décrochage de SGD en phase de raffinement
est la signature expérimentale de ce conditionnement.

\textbf{II $\to$ V (diagonal $\to$ quasi-Newton).} Ajoute des corrélations de
courbure au prix de $\sim$20 évaluations par pas. Sur ce problème, LBFGS
égale Adam sans le dépasser~: l'essentiel du conditionnement est
\emph{diagonal} (échelles par coordonnée), déjà capté par Adam~; la courbure
croisée n'apporte pas de gain mesurable à budget égal.

\subsection{Recommandations}

\begin{enumerate}
  \item \textbf{Pour une baseline de littérature}~: architecture I — solveur
  canonique, garanties, zéro réglage~; c'est la référence attendue par les
  relecteurs.
  \item \textbf{Pour la performance et l'extensibilité}~: architecture II —
  meilleure fidélité, zéros exacts, GPU, et la boucle d'entraînement se
  transpose telle quelle aux modèles non-linéaires (seul le décodeur change).
  \item \textbf{À reporter systématiquement ensemble}~: I \emph{et} II. Leur
  écart quantifie la part du solveur~; l'écart entre II et un modèle
  non-linéaire quantifie alors proprement la part de la classe de modèle.
  Confondre les deux écarts (comparer un non-linéaire optimisé par gradient à
  un linéaire résolu par ALS) surestime structurellement l'apport de la
  non-linéarité.
\end{enumerate}

% ══════════════════════════════════════════════════════════════════
\section{Détails d'implémentation}
% ══════════════════════════════════════════════════════════════════

\paragraph{Produits de mode.} La reconstruction (\ref{eq:tucker}) et ses
dégradations sont réalisées par contractions \texttt{einsum}
(\texttt{abc,ia,jb,kc->ijk})~; l'autodiff fournit exactement
(\ref{eq:gradG}) sans dérivation manuelle.

\paragraph{Double précision d'échelle.} Les entrées sont normalisées par
percentile~99 par source~; la sortie est dénormalisée par $s_H$ et écrêtée à
$[0,1]$.

\paragraph{Arrêt.} Budget maximal de $10^4$ itérations, arrêt anticipé si la
perte ne s'améliore plus de $10^{-7}$ pendant la patience~; l'état retenu est
celui de la \emph{meilleure perte} rencontrée (et non le dernier).

\paragraph{Suivi.} À chaque 10 itérations sont journalisés la perte couplée
et le taux de zéros exacts de $\T{G}$ — ce dernier étant le diagnostic de
santé du c\oe ur (une sparsité qui dérive vers 100\,\% signale un seuillage
pathologique).

\paragraph{Reproductibilité.} Mêmes graines, mêmes opérateurs
$\M{B}_h,\M{B}_w,\M{R}$ (construits depuis la PSF et la SRF \emph{fournies}
par le cadre d'évaluation), même initialisation
(\ref{eq:init-svd})--(\ref{eq:init-core}) pour les cinq architectures~: la
seule variable libre entre architectures est l'algorithme de descente.

% ══════════════════════════════════════════════════════════════════
\section{Conclusion}
% ══════════════════════════════════════════════════════════════════

Nous avons présenté cinq architectures de résolution d'un même modèle
linéaire de fusion — la décomposition de Tucker couplée à c\oe ur
parcimonieux — en explicitant l'intégralité des équations de mise à jour et
les conséquences de chaque choix algorithmique. Trois enseignements
principaux se dégagent~:

\begin{enumerate}
  \item \textbf{Le solveur fait partie du modèle empirique.} À paramétrisation
  identique, le passage de l'alternance exacte au gradient simultané
  préconditionné change les performances de plusieurs dB~: toute comparaison
  inter-modèles doit fixer le solveur ou reporter les deux.
  \item \textbf{La parcimonie exige un traitement proximal.} Seul le seuillage
  doux — dans FISTA comme après Adam — produit un support creux exact~; la
  pénalité sous-différentiée ne fait que rétrécir. Le découplage du seuil et
  du scheduler est nécessaire pour que $\beta$ garde un sens tout au long de
  l'entraînement.
  \item \textbf{Le cadre gradient est le pont vers le non-linéaire.} La perte
  couplée, les opérateurs physiques différentiables, l'étape proximale et la
  boucle d'entraînement sont réutilisables à l'identique lorsque le décodeur
  multilinéaire (\ref{eq:tucker}) est remplacé par un décodeur
  convolutionnel~: les architectures II--V constituent ainsi la base
  méthodologique commune des modèles non-linéaires du projet.
\end{enumerate}

\appendix

% ══════════════════════════════════════════════════════════════════
\section{Pseudo-code des deux schémas principaux}
% ══════════════════════════════════════════════════════════════════

\begin{figure}[h]
\centering
\fbox{\begin{minipage}{0.92\textwidth}
\textbf{Algorithme 1 — Architecture I : ALS + FISTA}
\begin{enumerate}\itemsep2pt
  \item Initialiser $(\M{D}_1,\M{D}_2,\M{D}_3,\T{G})$ par
        (\ref{eq:init-svd})--(\ref{eq:init-core}).
  \item \textbf{Pour} $i = 1,\dots,n_{\text{outer}}$ :
  \begin{enumerate}
    \item \textbf{Pour} $n = 1,2,3$ : $\M{D}_n \leftarrow$ solution de
          (\ref{eq:als-D})~; normaliser les colonnes (\ref{eq:normcols}).
    \item $\T{G} \leftarrow \T{G}/\norf{\T{G}}$~;
          calculer $L$ par l'équation (\ref{eq:lipschitz}).
    \item \textbf{Pour} $j = 1,\dots,n_{\text{FISTA}}$ :
          itération (\ref{eq:fista}) avec seuil $\beta/L$.
  \end{enumerate}
  \item \textbf{Retourner}
        $\T{S}=\T{G}\times_1\M{D}_1\times_2\M{D}_2\times_3\M{D}_3$.
\end{enumerate}
\end{minipage}}
\end{figure}

\begin{figure}[h]
\centering
\fbox{\begin{minipage}{0.92\textwidth}
\textbf{Algorithme 2 — Architectures II/IV/V : gradient + prox
(Adam / SGD / LBFGS)}
\begin{enumerate}\itemsep2pt
  \item Initialiser $\theta^{(0)}$ par
        (\ref{eq:init-svd})--(\ref{eq:init-core}).
  \item \textbf{Pour} $t = 1,\dots,T$ \textbf{ou} arrêt anticipé :
  \begin{enumerate}
    \item passe avant~: normaliser les colonnes \emph{dans le graphe},
          reconstruire $\T{S}$, dégrader, évaluer $f(\theta)$~;
    \item rétropropagation~; écrêtage $\lVert\nabla\rVert_2\le 10$~;
    \item pas de l'optimiseur (éq.~\ref{eq:adam}, \ref{eq:sgd} ou
          \ref{eq:lbfgs})~;
    \item $\T{G} \leftarrow \st_{\beta\eta_0}(\T{G})$
          \hfill (seuil découplé du scheduler)
    \item mémoriser $\theta$ si $f$ est la meilleure valeur rencontrée.
  \end{enumerate}
  \item \textbf{Retourner} $\T{S}(\theta^{\text{best}})$.
\end{enumerate}
\end{minipage}}
\end{figure}

\end{document}
